\chapter{A Premise}
Chapter \ref{ch:complexity} of this thesis has been written as a scientific paper and it is yet to be published.
The core of the paper has been written two years ago and looking at it now it appears limited in many of its thesis.

The possibility of an objective knowledge of reality is not yet refused and any insight given by \textcite{feyerabend1982} and \textcite{lyotard2005} are not incorporated.
What is there called complexity, I now refer to as a \textit{phenomenological approach}, in which the researcher's positioning is itself part of the knowledge produced.
The levels there pointed out can also be intended as subjective perspectives and the inability to model the world in a single unified way as a necessity rather than a shortcoming.

Hegemony and ideology did not find space in the original paper, but they are the elephant in the room: the role of cultural landscape and power relations in shaping what is thinkable (and acceptable) and what can be a subject of research is crucial in the production of knowledge.

%%

A broader critique of the paper (and generally of the methodology this thesis is rooted in) is the topic of chapter \ref{ch:against}.